\paragraph{\trVanishingCycle{}。}\label{sub:vanishing-cycles}
设 $\pi\colon E^{2n+2}\to S$ 是一个（照常\trExact{}的）\trLefschetzFibration。一条%
\textbf{\trVanishingPath}（vanishing path）是一条驶入\trCriticalPoint{}的\trEmbedding{}路径%
$\gamma\colon[0;1]\to S$，更确切地说，即满足 $\gamma^{-1}(\mathrm{Critv}(\pi))=\{1\}$ 的路径。对每条这样的路径都可以关联其\textbf{\trLefschetzThimble}（Lefschetz thimble）$\Delta_\gamma$，即 $E$ 中唯一的满足 $\pi(\Delta_\gamma)=\gamma([0;1])$、$\pi(\partial\Delta_\gamma)=\{\gamma(0)\}$ 的\trEmbedding\trLagrange{}$(n{+}1)$ 维球。其边界 $V_\gamma=\partial\Delta_\gamma$ 是 $E_{\gamma(0)}$ 中的一个\trLagrange\trSphere，称为 $\gamma$ 的\textbf{\trVanishingCycle}（vanishing cycle）。由于它界定 $E$ 中的一个\trLagrange\trDisk{}，任何\trVanishingCycle{}都自动是\trExact{}的（当然，这只在 $n=1$ 时才非平凡）。在继续往下之前，我们注意一个初等的事实：
\begin{lemma}\label{thm:critical-point-restriction}
对任何\trLagrange\trSubmanifold{}$F\subset E$，点 $x\in F$ 是 $\pi$ 的\trCriticalPoint{}当且仅当它是 $\pi|F$ 的\trCriticalPoint{}。
\end{lemma}
\noqed
就具有上述性质的\trDisk{}$\Delta_\gamma$ 而言，函数 $p=\gamma^{-1}\circ\pi|_{\Delta_\gamma}\colon\Delta_\gamma\to[0;1]$ 必在内部某处取得最大值，而由\cref{thm:critical-point-restriction}，这个点必是 $E_{\gamma(1)}$ 中唯一的\trSingularPoint{}$x_1$。一个类似的论证表明，$p$ 在该点的\trHessian{}是\trNondegenerate{}的，故 $p$ 是一个只有单个\trCriticalPoint{}的\trMorseFunction{}。在给 $\Delta_\gamma$ 装备一个\trMetric{}——它在\trCriticalPoint{}附近的局部\trMorse{}坐标中是标准的——之后，就可以用 $\nabla p$ 的\trFlow{}来提供一个\trDiffeomorphism
\begin{equation}\label{eq:thimble-ball-diffeomorphism}
B=B(T_{x_1}\Delta_\gamma)\longrightarrow\Delta_\gamma
\end{equation}
它把 $x_1$ 处\trTangentSpace{}中的单位球映到整个\trLefschetzThimble{}。特别地，限制到边界上便给出一个\trFrame{}$v_\gamma\colon S=\partial B\to V_\gamma$，它在相差\trIsotopy{}的意义下不依赖于\trMetric{}的选取。我们把 $v_\gamma$ 取作\trVanishingCycle{}之结构的一部分，因此后者今后总被视为一个带\trFrame{}的\trExact\trLagrange\trSphere。
我们已经从\trLefschetzThimble{}的定义性质推出了若干结论，但它的存在性与唯一性仍有待验证；我们现在就来处理这件事，从唯一性开始。下面是一个有用的事实：
\begin{lemma}\label{thm:fibered-lagrangian-criterion}
设 $\beta$ 是底 $S$ 的\trRegular{}部分中的一条\trEmbedding{}路径，而 $F\subset E$ 是一个在 $\beta$ 上\trFibration{}的\trSubmanifold{}，使得每个 $F_{\beta(s)}\subset E_{\beta(s)}$ 都是\trLagrange{}的。那么整个 $F$ 是\trLagrange{}的，当且仅当沿 $\beta$ 的任何一段作的\trParallelTransport{}都把诸 $F_{\beta(s)}$ 彼此相互映射。
\end{lemma}
\noqed
对\trLefschetzThimble{}而言，这意味着沿 $\gamma$ 的\trParallelTransport{}把诸\trLevelSet{}$\Delta_\gamma\cap E_{\gamma(t)}=p^{-1}(t)$（$0\le t<1$）彼此相互映射。当 $t\to 1$ 时，这些\trLevelSet{}缩为 $\{x_1\}$。因此，$\Delta_\gamma$ 必然包含于集合
\begin{equation*}
\Delta_\gamma'=\{x\in E_{\gamma(t_0)},\,0\le t_0<1:\lim_{t_1\to 1}h_{\gamma|[t_0;t_1]}(x)=x_1\}\cup\{x_1\}.
\end{equation*}
$\Delta_\gamma'$ 的结构可能不是一眼可见的，因为\trParallelTransport{}\trVectorField{}在 $x_1$ 处有一个极点。为绕开这个困难，选一个\trSmooth{}函数 $h\colon S\to\RR$，使 $h(\gamma(t))=0$，并且使 $(\gamma'(t),\nabla h_{\gamma(t)})$ 对任何 $t$ 都构成切向量的一个正\trOrientation{}基。令 $X_H$ 为 $H=h\circ\pi$ 的\trHamiltonianVectorField{}。容易看出，在 $\pi^{-1}(\gamma([0;1)))$ 上，它指向与\trParallelTransport{}\trVectorField{}相同的方向，同时在 $x_1$ 处有一个\trStationary{}点。于是 $\Delta_\gamma'$ 包含于
\begin{equation}\label{eq:thimble-stable-manifold}
\Delta_\gamma''=\{x\in E_{\gamma(t)},\,0\le t\le 1:\lim_{s\to\infty}\varphi_H^s(x)=x_1\}.
\end{equation}
由于\trStationary{}点 $x_1$ 是双曲的，可以应用\trStableManifold{}定理来证明 $\Delta_\gamma''$ 是一个\trEmbedding{}的 $(n{+}1)$ 维球，其边界为\trSphere{}$\Delta_\gamma''\cap E_{\gamma(0)}$。由于按假设 $\Delta_\gamma$ 也是如此，包含关系 $\Delta_\gamma\subset\Delta_\gamma'\subset\Delta_\gamma''$ 必为等式，这就证明了\trLefschetzThimble{}的唯一性。存在性则这样处理：注意到 \eqref{eq:thimble-stable-manifold} 是\trLagrange{}的（因为辛\trFlow{}$\varphi_H$ 在其上是压缩的），并且具备其余全部想要的性质。
\begin{remark}\label{thm:vanishing-cycle-piece-transport}
作为上述讨论的一个副产品，我们看到任一段 $\gamma|[t_0;1]$ 的\trVanishingCycle{}与整条路径的\trVanishingCycle{}通过\trParallelTransport{}相联系：
\begin{equation}\label{eq:vanishing-cycle-transport}
V_\gamma=h_{\gamma|[0;t_0]}^{-1}\bigl(V_{\gamma|[t_0;1]}\bigr).
\end{equation}
\end{remark}
这一观察可以用来为未必\trEmbedding{}的路径 $\gamma$ 定义\trLefschetzThimble{}与\trVanishingCycle{}——只要它仍满足 $\gamma^{-1}(\mathrm{Critv}(\pi))=\{1\}$ 与 $\gamma'(1)\neq 0$。也就是说，取 $t_0$ 充分接近 $1$，则 $\gamma|[t_0;1]$ 将是一条\trEmbedding{}路径，从而可以用\trParallelTransport{}把 $\Delta_{\gamma|[t_0;1]}$ 延拓出去。这样定义的诸\trLefschetzThimble{}是\trImmersion{}的、但未必\trEmbedding{}的\trLagrange{}球；不过它们的边界——恰好就是 \eqref{eq:vanishing-cycle-transport}——将具有与\trEmbedding{}情形相同的性质。
\begin{example}\label{thm:toy-model-thimble}
取 \cref{thm:basic-local-model} 中定义的玩具模型 $\pi\colon E\to S$。这并不完全是一个\trLefschetzFibration{}，但所缺的条件（缺少\trHolomorphic\trConvex{}性）与本处的目的无关。唯一的\trCriticalValue{}是 $0$，而对任何\trVanishingPath{}$\gamma$，\trLefschetzThimble{}都可以显式确定：
\begin{equation}\label{eq:toy-model-thimble}
\Delta_\gamma=\bigcup_{0\le t\le 1}\sqrt{\gamma(t)}\,S^n;
\end{equation}
这里 $\sqrt{\bar{z}}\,S^n=\{x\in\CC^{n+1}:x=\pm\sqrt{\bar{z}}\,y\}$，其中 $y$ 是 $S^n\subset\RR^{n+1}$ 中的某个元素。为看出 \eqref{eq:toy-model-thimble} 成立，只需用 \cref{thm:basic-local-model} 中的函数 $k$——它在\trParallelTransport{}下不变——并注意到 $k^{-1}(0)$ 恰好是所有 $\sqrt{\bar{z}}\,S^n$（$z\in S$）的并。
\end{example}
